\ssrno{ВВЕДЕНИЕ}

Распространение синтезированного контента, созданного с использованием технологий машинного обучения, приводит к увеличению числа правонарушений и экономическому ущербу~\cite{spb-uni}. Количество дел, связанных с применением искусственного интеллекта, в российских судах за последние пять лет возросло более чем на 40\%. Визуальный анализ не позволяет достоверно идентифицировать значительную часть качественно исполненных поддельных изображений. Указанные обстоятельства обусловливают необходимость разработки автоматизированных методов обнаружения изображений-имитаторов.

Целью данной выпускной квалификационной работы является разработка метода обнаружения признаков генерации на изображении с использованием нейронной сети.

Для достижения поставленной цели необходимо выполнить следующие задачи:
\begin{itemize}
	\item проанализировать предметную область;
	\item проанализировать известные нейросетевые решения;
	\item провести классификацию и сравнительный анализ рассмотренных решений;
	\item формализовать задачу в виде функциональной модели;
	\item спроектировать и разработать программное обеспечение для обнаружения признаков генерации на изображении;
	\item провести исследование работоспособности программной реализации предложенного метода.
\end{itemize}